\subsection{Certified removal of coordinate factors}
The finite fibres can certify further degree reductions before the
bounded-excess algorithm is used. All powers of a coordinate may be cancelled uniformly from equations
in the prime ideal $I_X$ whenever that coordinate does not vanish
identically on $X$. The unknown coordinates that do
vanish span a space of dimension at most $c-1$; the following algorithm
finds certificates without first computing $I_X$.

Fix a deterministic preprocessing convention on the displayed generating
list of $K$. For a set
$H\subseteq\{0,\ldots,N\}$, process each nonzero row by dividing out its
full monomial content in coordinates outside $H$. Then apply the
rank-one digit operation of Lemma~\ref{sci:lem:deflation}, choosing the
largest admissible $r\ge1$. Then apply a minimum-degree recognized
replacement from Theorems~\ref{sp:thm:linearized-preprocessing},
\ref{sp:thm:affine-preprocessing}, and~\ref{sp:thm:zero-sum-preprocessing}, breaking ties by a fixed
ordering. Normalize each output row by its first nonzero coefficient
and repeat on the resulting list until no operation applies. Let
$\mathcal P_H(K)$ denote the ideal of the resulting list. Here the digit
operation writes $f=AB^{p^r}$ with all exponents of $A$ less than $p^r$
and $B$ nonconstant, and replaces it by $AB$. It is recognized by the
rank-one coefficient matrix of low and high exponent digits. Zero rows
are skipped. This convention defines the processed ideal even when the
speculated coordinate divisions are not valid on $X$.
Discard any speculative branch that produces a nonzero constant.

\begin{theorem}[Certified coordinate division before saturation recovery]
\label{sci:thm:certified-preprocessing}
Fix a prime $p$, a codimension $c\ge2$, and $E\ge0$. Let a homogeneous
sparse list with binary exponents generate
$K\subset\Fp[x_0,\ldots,x_N]$, with
$K^{\mathrm{sat}}=I_X$, where $X$ is a smooth geometrically integral
anticanonical global complete intersection of dimension $N-c\ge1$.
Set
\[
 H_X=\{j:x_j\in I_X\},\qquad K^\circ=\mathcal P_{H_X}(K).
\]
Suppose that, after a finite field extension, $K^\circ$ contains a
regular $c$-sequence of degree sum at most $N+1+E$. Neither $H_X$ nor
that sequence is part of the input. For every $a\ge1$, a randomized
algorithm of bit cost polynomial in the written input length and $a$
returns the coherent Frobenius scalar or a failure symbol, with failure
probability at most $2^{-a}$ and no incorrect outputs on the promise.
Repetition has expected polynomial cost. Its only random choices are
projection centres. Every coordinate division used for the final
scalar is certified from the original equations on a deterministic
grid of at most $N+1$ affine fibres for each tested tuple and centre.
\end{theorem}
\begin{proof}
Write $I=I_X$, $n=N-c$, and $R=N+1+E$. The space $I_1$ has dimension
at most $c-1$. Otherwise all $c$ minimal complete-intersection
generators would be linear, contradicting their degree sum
$N+1\ge c+2$. Consequently $|H_X|\le c-1$, and there are
$\sum_{j=0}^{c-1}\binom{N+1}{j}$ possible sets to examine.
Every coordinate occurs in the original list. Indeed, if $x_j$ were
absent, multiplication by it would be injective on $S/K$, so saturation
would give $K=I$. The coordinate point $e_j$ would then be a cone vertex
whose tangent space has dimension $N-\dim I_1\ge N-c+1$, contradicting
smoothness of the $(N-c)$-dimensional variety. Thus $N+1$ itself is
bounded by the written sparse input length, even with sparse coordinate
labels.
For the true set, every division is valid: if $x_j\notin I$ and
$x_j^b u\in I$, primality gives $u\in I$. Digit reduction is valid by
radicality, and the recognized prime-field relations are valid by
geometric integrality. Every old row belongs to the ideal generated by its replacement rows, so the
ideal sandwich $K\subseteq K^\circ\subseteq I$ proves
$(K^\circ)^{\mathrm{sat}}=I$.

We first establish the certificate used even on a wrong branch. Let
$k/\Fp$ be a finite extension, let $t_0$ be any linear form over $k$,
and let $B$ be a finite-dimensional quotient of
\[
 S_k/(K_k,t_0-1).
\]
If the image of $x_j$ in $B$ is nonzero, then $x_j\notin I$. Indeed,
$x_j\in K^{\mathrm{sat}}$ would give
$\mathfrak m^b x_j\subset K$ for some $b$, and hence
$t_0^b x_j\in K_k$; the image would vanish after setting $t_0=1$.
The certificate is a nonzero algebra element. It need not be a unit,
and nilpotents are not discarded.

For each possible $H$, form the speculative list $\mathcal P_H(K)$.
Discard a branch containing a nonzero constant. Run the explicit candidate construction of
Lemma~\ref{sci:lem:profile-candidates} on all profiles of sum between
$N+1$ and $R$. For each candidate tuple, sample a centre and apply the
exact degree-bounded Macaulay certificate described in
Proposition~\ref{sci:prop:random-designated}.
A successful certificate proves that the tuple is regular and finite
over its projection base, whether or not the speculative list has the
correct saturation. Complete its coordinates as $x=Uz+Vt$, where
$t=(t_0,\ldots,t_n)$. For
\[
 \Gamma_1=\{0,e_1,\ldots,e_n\}\subset k^n
\]
with $e_i$ the standard coordinate vectors,
form the full finite algebra
\[
 A_b=k[z_1,\ldots,z_c]/
       (f_i(Uz+V(1,b)):1\le i\le c).
\]
Its dimension is $\Delta=\prod_i\deg f_i\le R^c$.
Evaluate every original input row in $A_b$ by binary powering and
quotient by the ideal of those values, obtaining $B_b$. This is exact
finite-dimensional linear algebra: if $v_1,\ldots,v_\Delta$ is an
algebra basis, the vectors $v_j u_i$ span the ideal generated by the
values of the original rows. Compute the image of every coordinate in
this quotient. Let $C$ collect the indices having a nonzero image in
at least one of these algebras, over all branches examined.
Each $B_b$ is a quotient of $S_k/(K_k,t_0-1)$, so every index in $C$
is certified safe even if its branch was wrong.

For completeness, take the true branch and one of its regular
bounded-excess candidates with a successful centre. Its ideal $J$ is
contained in $I$. The graded module $M=S_k/I_k$ is finite free over
$k[t_0,\ldots,t_n]$, with a homogeneous basis of nonnegative degrees,
by the finite-free argument of Lemma~\ref{sci:lem:fibres}.
On $t_0=1$, the original ideals $K_k$ and $I_k$ agree: a power of the
irrelevant ideal annihilates their quotient, and $t_0$ is invertible.
Therefore the certificate algebra is exactly
\[
 B_b=M\otimes_{k[t_0,\ldots,t_n]} k_{(1,b)}.
\]
In a homogeneous free basis, a degree-one coordinate has coefficient
polynomials of degrees at most one, because every basis degree is
nonnegative. If $x_j\notin I$, at least one coefficient is nonzero.
Dehomogenization at $t_0=1$ is injective on a fixed homogeneous degree,
and evaluation at $0,e_1,\ldots,e_n$ detects every nonzero affine linear
polynomial over every field. Thus one of these fibres certifies $j$.
On this event $C$ is exactly the complement of $H_X$.

Process the original list again, dividing only by coordinates in $C$
and applying digit and prime-field relation reductions with the fixed
convention. All these
operations now preserve saturation, independently of whether $C$ is
complete. Run Theorem~\ref{sci:thm:random-excess} on this safe list.
Every successful output is correct. Once the true branch has a good
centre, its list is exactly $K^\circ$, so the assumed bounded-excess
promise applies to this final run.

Assign failure budgets $2^{-a-1}$ to certificate discovery and to the
final scalar calculation. Only one regular candidate on the true branch
is needed for discovery; wrong branches cannot produce a false
certificate. The centre resultant has degree at most $cR^c$, so the
same finite-field sampling and repetition bound used in
Theorem~\ref{sci:thm:random-excess} gives its required failure budget.
A union bound between the two stages proves the stated probability.
Both stages return a distinguished failure when their bounded trials
are exhausted; no unbounded search is hidden in this assertion.

It remains to check the binary-degree cost of preprocessing.
Theorem~\ref{sp:thm:zero-sum-preprocessing} bounds both total support
and the number of steps polynomially through the entire interleaving,
including affine row splitting and zero-sum support increases.
These combinatorial bounds apply on speculative branches up to rejection
as well: all surviving rows are homogeneous of positive degree,
regardless of membership in $I$.
Its token count uses the stronger square-root degree decrease at
every potentially support-increasing step. Full eligible monomial
content is removed in maximal batches between degree reductions.
Testing every possible $r$ uses only
binary exponent arithmetic and sparse rank-one tests. The number of
branches, profiles, algebra dimensions, coordinate tests and fibre
points is polynomial for fixed $c,E$. Original large exponents occur
only in binary powers of finite-algebra elements. Finite-field
construction and scalar descent have the already established polynomial
bit bounds. This completes the proof.
\end{proof}

\begin{corollary}[A deterministic processed linear prefix]
\label{sci:cor:certified-linear}
Under the hypotheses of Theorem~\ref{sci:thm:certified-preprocessing},
if $c\ge2$ and $\dim(K^\circ)_1=c-1$, the scalar is computable
deterministically in polynomial bit time.
\end{corollary}
\begin{proof}
Enumerate the same possible protected-coordinate sets. On each branch
whose degree-one part has dimension $c-1$, run the bounded-loop
centre construction of Theorem~\ref{sci:thm:linear-prefix}, testing
its exact centre certificate and returning failure if any required
step fails. The true branch has a regular bounded-excess tuple and
succeeds. Its degree-one fibre grid certifies every safe coordinate.
All certificates from other branches remain sound by the preceding
proof. Reprocessing with their union therefore gives $K^\circ$.
Apply the deterministic linear-prefix algorithm to that list. Each
stage has a fixed polynomial bound even on a branch failing its
promises; only verified centres and coordinate certificates are used.
\end{proof}

\begin{example}[Unbounded excess before certified division]
\label{sci:ex:certified-fermat}
Fix $p$ and choose arbitrarily large $N\ge3$ with $p\nmid N$ and
$N\equiv1\pmod{p-1}$, for example $N=1+p(p-1)k$, $k\ge1$.
In $\Fp[w,x_1,\ldots,x_N]$ put
\[
 F=\sum_{i=1}^N x_i^N,\quad I=(w,F),\quad
 z=(w,x_1,\ldots,x_N),\quad
 K=(z_j^{N-1}w,z_j^{N-1}F:0\le j\le N).
\]
Then $K^{\mathrm{sat}}=I$, its anticanonical variety is smooth and
geometrically integral, and the least complete-intersection excess
after digit reduction alone is exactly $N-1$. Certified coordinate
division gives excess zero and the deterministic case of
Corollary~\ref{sci:cor:certified-linear}.
\end{example}
\begin{proof}
On every projective coordinate chart a cover factor is a unit, so the
saturation is $I$. The nonzero gradient entries $Nx_i^{N-1}$ prove
smoothness of the degree-$N$ hypersurface in $w=0$. Two distinct
positive-degree factors would intersect in projective space and give
a singular point; repeated factors are also impossible. Thus it is
geometrically integral. The degrees $1,N$ sum to $N+1$.

The row $w^N$ digit-reduces to $w$: a positive monomial exponent keeps
its residue modulo $p-1$, and a terminal positive exponent lies between
one and $p-1$. Every $F$-row is unchanged by a nontrivial rank-one digit
test. To see this, fix $P=p^r$ and write $N=dP+b$ when $P\le N$.
Then $d\ge1$ and $1\le b<P$, since $p\nmid N$.
For $x_j^{N-1}F$, the terms with $i\ne j$ have distinct low digits
$x_j^{b-1}x_i^b$ and distinct high digits $x_j^d x_i^d$.
No other term shares these rows or columns of the coefficient matrix,
so two such terms give rank at least two. They exist because $N\ge3$.
The $w$-cover row has the same property with its $N$ distinct
$x_i$ terms. If $N<P\le2N-1$, the internal row has one term with high
digit $x_j$ and the remaining terms have high digit one; their low
supports are disjoint, so the rank is again two. The $w$-cover row
has constant high factor in this range. For $P>2N-1$ all high factors
are constant. These are all possible digit tests.

After digit reduction the ideal therefore contains $w$, and every
remaining row outside $(w)$ has degree $2N-1$. Every homogeneous
element of smaller degree lies in $(w)$. A regular pair consequently
has degree sum at least $1+(2N-1)=2N$, attained by
$(w,x_1^{N-1}F)$. Its excess is $N-1$. On the other hand,
$H_X=\{0\}$, where index $0$ denotes the coordinate $w$, and safe
division removes $x_i^{N-1}$ from an internal
$F$-row. The fully processed ideal is $(w,F)$, of excess zero and
with a one-dimensional degree-one part, as asserted.
\end{proof}
